\chapter{Agentix Islam: Architecting the Zero-Hallucination Framework}
\label{chap:Agentix_Islam:_Architecting_the_Zero-Hallucination_Framework}

\section{Technical Implementation of the Agentic Islamic Fatwa (RAG) Framework}

This chapter details the exact methodology, architecture, and logic used to construct the custom Retrieval-Augmented Generation (RAG) system for the Islamic Fatwa agent (corresponding to the \texttt{R6-Agentix Islam (RAG)} project).

\subsection{Data Ingestion and Structured Parsing}
The foundational step involved digitizing and parsing classical Islamic jurisprudence texts into a machine-readable format while strictly preserving scholarly integrity.

\begin{itemize}
    \item \textbf{Batch Processing with LLMs:} The system uses the \textbf{Gemini 2.5 Pro} model via the Batch API to process raw manuscript pages at scale.
    \item \textbf{JSON Schema Enforcement:} To ensure structural consistency, a strict JSON schema was provided in the system prompt. The model was instructed to extract content page-by-page into a strict dictionary containing:
    \begin{itemize}
        \item \texttt{page\_content}: the main text of the page.
        \item \texttt{page\_number}: the specific, verifiable page number to ensure accurate citation later.
        \item \texttt{notes}: a specialized sub-object designed to uniquely capture classical text annotations, separating \texttt{footnote}s (Footnotes) from \texttt{marginalia} (Marginalia).
    \end{itemize}
\end{itemize}


\subsection{Hierarchical Metadata and Table of Contents (TOC) Engineering}
An intelligent RAG system requires semantic understanding of the book's structure, not just isolated flat text strings.

\begin{itemize}
    \item \textbf{TOC Flattening:} The books' structural Tables of Contents were mapped into a hierarchical JSON array. Each section, chapter, or topic (e.g., ``Introduction to the New Edition'') was mapped with a \texttt{title}, \texttt{startPage}, \texttt{endPage}, and its nesting \texttt{level} or \texttt{parentId} (pathIds and pathTitles).
    \item \textbf{Markdown Synthesis:} The extracted raw JSON pages were subsequently stitched together to form a clean, continuous Markdown structure. Page numbers and separated footnotes were interleaved at designated points, maintaining the text's academic validity while making it parsable for chunking algorithms.
\end{itemize}

\subsection{Vectorization and Database Construction}
Before being queried, the structured text was transformed into embeddings designed for high-precision semantic search.

\begin{itemize}
    \item \textbf{Intelligent Text Splitting:} The compiled Markdown documents were chunked using LangChain's \texttt{RecursiveCharacterTextSplitter}. Because Islamic jurisprudence requires high context retention, the chunks were set to 1000 characters with a safety overlap of 200 characters to prevent cutting off crucial conditions or \textit{Ahadith}.
    \item \textbf{Metadata Tagging:} Every individual chunk was injected with localized metadata referencing its origin: \texttt{book\_id}, \texttt{book\_name}, \texttt{book\_part\_number} (volume), and the exact \texttt{page\_number}.
    \item \textbf{Embeddings and Storage:} Text chunks were vectorized using Google's \texttt{gemini-embedding-001} model and stored immutably into a cloud-hosted \textbf{Chroma DB} (with an architectural fallback to local FAISS storage).
    \item \textbf{Two-Stage Retrieval Mechanism (MMR + Reranking):} To solve the pervasive issue of irrelevant fatwa generation, the retrieval strategy utilizes a rigorous pipeline:
    \begin{enumerate}
        \item \textbf{Initial Retrieval:} Maximum Marginal Relevance (MMR) fetches the top \textit{k=5} diverse but highly relevant chunks.
        \item \textbf{Contextual Compression:} A \texttt{CohereRerank} (multilingual v3.0) model is layered on top to re-score and compress the returned nodes, pushing the most semantically relevant theological rulings to the top of the context window.
    \end{enumerate}
\end{itemize}

\subsection{Agentic Orchestration and Backend Logic (FastAPI)}
The central intelligence of the framework is a dynamic orchestrator built with FastAPI, deploying \textbf{Gemini 3 Flash Preview} (and Gemini 2.5 Flash). It functions autonomously as an ``Islamic Scholar'' agent.

\begin{itemize}
    \item \textbf{Strict RAG Guardrails:} The LLM is given a definitive system instruction commanding it to act strictly as an Islamic scholar. It is forced to answer \textbf{only} from the retrieved context. If the ruling does not exist, it must declare its inability to answer (``Information not available in the provided sources'') instead of hallucinating. It must answer in Arabic or English and maintain scholarly etiquette.
    \item \textbf{Function Calling (The Agentic Component):} Rather than blindly performing similarity searches on every query, the LLM is equipped with specialized tools (functions):
    \begin{itemize}
        \item \texttt{query\_vector\_store}: Triggered automatically for general theological queries across the whole database.
        \item \texttt{find\_page\_scopes} and \texttt{get\_pages\_in\_range}: Triggered if the user asks for a deep dive, or specifies a book part and page. The LLM navigates the previously defined TOC JSON tree to fetch specific, contiguous pages.
    \end{itemize}
    \item \textbf{Verification and Output Generation:}
    \begin{itemize}
        \item The LLM enforces strict citation mechanisms natively structured from the chunk metadata (e.g., ``(Page 584, Vol 1)'').
        \item It generates its final response in a strictly validated JSON format ensuring three keys: \texttt{answer} (the theological ruling in markdown), \texttt{relevant\_questions} (generating three contextually aware follow-up discussion points to keep the user engaged in a learning loop), and \texttt{resources} (an array tracking all accessed page numbers and volumes).
        \item A continuous database layer (\texttt{TocManager}) intercepts the LLM's tool usages, storing token counts, search trajectories, and conversation history asynchronously using FastAPI's \texttt{BackgroundTasks}.
    \end{itemize}
\end{itemize}
